\section{从递归到 Attention}

\subsection{RNN 时代的主流做法}

在 Transformer 之前，NLP 的主流方案是用双向 LSTM 编码输入，再用解码器逐步生成输出。这个范式虽然有效，但它有两个明显问题：难以并行，且长距离依赖的梯度传播路径太长。

\begin{figure}[H]
\centering
\includegraphics[width=0.95\textwidth]{slides-images/slide-016.jpg}
\caption{2016 年前后 NLP 的主流范式：双向 LSTM 编码器 + LSTM 解码器\protect\footnotemark}
\end{figure}
\footnotetext{Slides 第17页。}

\begin{figure}[H]
\centering
\includegraphics[width=0.95\textwidth]{slides-images/slide-023.jpg}
\caption{RNN 与 attention 的计算依赖对比：recurrence 的串行依赖限制了并行化\protect\footnotemark}
\end{figure}
\footnotetext{Slides 第24页。}

\begin{importantbox}{RNN 的两个结构性瓶颈}
\begin{enumerate}
    \item \textbf{并行性差}：下一个状态必须依赖上一个状态，训练时很难充分利用硬件
    \item \textbf{交互距离长}：两个相距很远的词要通过许多时间步才能彼此影响，信息和梯度都更难传播
\end{enumerate}
\end{importantbox}

\subsection{Attention 作为过渡}

上一讲已经看到，attention 可以让 decoder 直接访问 encoder 中任意位置的信息。Transformer 进一步把这个思想推广到编码器和解码器内部，变成 \textbf{self-attention}。

\begin{figure}[H]
\centering
\includegraphics[width=0.95\textwidth]{slides-images/slide-022.jpg}
\caption{自注意力的高层结构：每个词都可以和序列中的其他词交互\protect\footnotemark}
\end{figure}
\footnotetext{Slides 第23页。}

\begin{figure}[H]
\centering
\includegraphics[width=0.95\textwidth]{slides-images/slide-024.jpg}
\caption{Transformer 的核心优势：不会随着序列长度增加而增加不可并行步骤，任意两个词之间的最大交互距离是 $O(1)$\protect\footnotemark}
\end{figure}
\footnotetext{Slides 第25页。}

\subsection{本章小结}

从 RNN 到 attention 的关键转变是：不再让信息只能沿着时间轴一步一步流动，而是允许每个位置直接检索整个序列。Transformer 把这个思想推到极致，于是计算图也变得更适合并行硬件。

%% ========== 理解 Transformer ========== 

